% ============================================================
%  Chapter 3 — NEXUS Architecture
%  STATUS: first draft (2026-05-21).
%  Source: DESIGN_DECISIONS.md (DD-V5-001 .. DD-V5-022),
%          docs/NARRATIVE_ARC_v5.md, docs/AGENT_CONSTRAINT_DERIVATION.md,
%          src/ implementation.
% ============================================================
\chapter{The NEXUS Architecture}
\label{ch:architecture}

% ------------------------------------------------------------
\section{Introduction}
\label{sec:arch-intro}

Manufacturing Operations Management (MOM) sits at Level~3 of the ISA-95
functional hierarchy, between the business-planning systems of Level~4
and the process-control systems of Levels~2 and below
\cite{isa95part1,sprock2019isa95}. Level~3 is the layer at which the
abstract production targets handed down by enterprise planning are
turned into concrete, time-bounded decisions on the plant floor: which
operation runs on which machine, when a machine is taken offline for
maintenance, whether a batch is held for quality inspection, and how far
ahead spare parts are ordered. ISA-95 decomposes this layer into four
operations-management functions --- production, maintenance, quality,
and inventory (here treated as supply) --- and models each as a
separate activity model with its own information flows and its own
performance criteria \cite{isa95part1}.

The decomposition is organisationally convenient but it is not
decisionally clean. The four functions are \emph{coupled}: a decision
that is locally optimal for one function routinely degrades another. A
maintenance window reclaims machine health but forgoes production
hours. A quality hold suppresses defect-driven rework but delays
throughput. A conservative supply preorder shortens stockout exposure
but raises carrying cost. A faster production rate lifts short-term
output but accelerates tool wear and, with it, both the breakdown
hazard and the defect rate. None of these trade-offs is visible inside
the activity model of a single function; each becomes visible only when
the four functions are considered jointly. This is the
\emph{cross-functional coordination gap} that this thesis addresses.

In most installed plant-floor systems the gap is left open. Each
function runs its own scheduler or its own rule set, optimising its own
key performance indicator, and the interactions between functions are
absorbed informally --- by shift supervisors, by escalation meetings,
or simply by whichever function happens to act first. The result is a
\emph{siloed} mode of operation in which the cross-functional
trade-offs are never explicitly represented, let alone resolved against
a shared objective. The coordination gap is not a modelling gap: a
sufficiently detailed centralised model could, in principle, optimise
the joint problem. It is a coordination gap, because the conditions
under which such a centralised model could be built --- full visibility
of every function's cost structure, machine state, defect history and
supplier reliability --- do not hold in practice. The functions are
organisationally distinct, their information systems are heterogeneous,
and the data each would have to surrender to a central optimiser is
exactly the data each treats as its own (Section~\ref{sec:arch-privacy}
returns to this point).

This chapter presents \NEXUS{} --- Negotiation-driven EXecution and
Unified Synchronization --- a reference architecture for closing the
cross-functional coordination gap at ISA-95 Level~3 without violating
the informational and organisational boundaries that make centralised
optimisation infeasible. \NEXUS{} is organised as four layers. A
\emph{Worker} layer hosts four Belief--Desire--Intention (BDI) agents,
one per ISA-95 Level~3 function. A \emph{Marketplace} layer hosts a
multi-issue negotiation mediator that aggregates the agents'
utility-annotated proposals and resolves each cross-functional conflict
into a single composite contract. A \emph{Fence} layer audits every
contract against a suite of pre-declared safety rules and classifies it
into a three-tier severity scheme. A \emph{School} layer --- described
here but deferred as future work --- provides the architectural slot
for continual adaptation of the agents' preferences from operational
history.

Two design commitments run through the whole architecture and are
worth stating at the outset.

\begin{description}[leftmargin=1.4em,style=nextline]
  \item[Composition over invention.] The contribution of \NEXUS{} is
    the \emph{composition} of established components into a coherent
    reference architecture for a specific problem --- not the invention
    of a new negotiation mechanism or a new learning algorithm. The
    negotiation mechanism in the Marketplace layer is the multi-issue
    multi-lateral protocol of Ito, Hattori and Klein
    \cite{ito2007multiissue}; the agent model is the BDI architecture
    of Rao and Georgeff \cite{rao1995bdi}; the reference-architecture
    framing follows the lineage of PROSA \cite{vanbrussel1998prosa},
    ADACOR \cite{leitao2007holonic} and HAPBA \cite{leitao2015hapba}.
    \NEXUS{} chooses, adapts and integrates; it does not reinvent. This
    commitment is deliberate and is defended throughout the chapter:
    a reference architecture earns its keep by being a sound,
    reproducible composition, not by maximising novelty in any single
    component.
  \item[Bounded information disclosure.] No component of \NEXUS{} ever
    requires a function to publish its internal cost model, its private
    state, or its objective weights. Agents disclose only
    \emph{bid summaries} --- a bounded number of regions of the joint
    decision space, each annotated with a single scalar utility. This
    is the structural property that distinguishes \NEXUS{} from a
    centralised optimiser and from distributed constraint optimisation,
    and it is treated as a first-class architectural requirement rather
    than an afterthought.
\end{description}

The remainder of the chapter is organised as follows.
Section~\ref{sec:arch-overview} presents the four-layer architecture as
a whole, fixes terminology, and traces the information flow of a single
coordination event. Section~\ref{sec:arch-worker} describes the Worker
layer and the four BDI agents, including the mono-dimensional output
discipline that keeps each agent within its ISA-95 domain.
Section~\ref{sec:arch-marketplace} describes the Marketplace layer: the
composite contract space, the bid representation, and the mediator's
admissible-heuristic search for the welfare-maximising joint contract.
Section~\ref{sec:arch-fence} describes the Fence layer and its
three-tier safety classification. Section~\ref{sec:arch-school}
describes the School layer at the level of detail needed to make it a
coherent future-work proposal. Section~\ref{sec:arch-isa95} consolidates
the mapping of the architecture onto the ISA-95 Level~3 functional model
and states the privacy-preservation property precisely.
Section~\ref{sec:arch-impl} closes with an overview of the
implementation architecture and forward references to
Chapter~\ref{ch:implementation}.

% ------------------------------------------------------------
\section{Architectural Overview}
\label{sec:arch-overview}

\subsection{The four layers}
\label{subsec:four-layers}

\NEXUS{} is a layered architecture. Each layer is a set of components
with a single, well-defined responsibility, and the layers communicate
only through narrow, explicitly typed interfaces. The layering is not
merely a presentational device: it is the mechanism by which the two
design commitments of Section~\ref{sec:arch-intro} are enforced. Because
the Worker layer can reach the Marketplace layer only by publishing bid
summaries, no agent can leak its cost model downstream even if it
wished to; because the Fence layer consumes only the Marketplace's
audit log, safety governance cannot couple itself into the negotiation
protocol. The interface discipline turns the design commitments into
structural invariants.

Figure~\ref{fig:arch-layers} shows the four layers and the principal
information flows between them. Table~\ref{tab:layers} summarises the
responsibility, the principal components, and the inbound and outbound
interface of each layer.

\nexusfigure{fig:arch-layers}{The four-layer \NEXUS{} reference
architecture. Solid arrows carry information exchanged during a
coordination event; the dashed boundary of the School layer marks it as
an architectural slot deferred to future work.}{a stacked
four-layer block diagram. Top to bottom: Worker layer (four BDI agent
boxes PSA, PMA, QCA, SCA); Marketplace layer (Bid Bus, Mediator,
Contract Record); Fence layer (three-tier classifier); School layer
(drawn with a dashed border). Upward arrows from Workers to Marketplace
labelled ``bid summaries''; a downward arrow from Marketplace to Workers
labelled ``composite contract''; an arrow from the Contract Record into
the Fence classifier labelled ``audit log''; dashed arrows from School
back to the Workers labelled ``adapted weights (future work)''.}

\begin{table}[ht]
\centering
\caption{The four layers of the \NEXUS{} architecture.}
\label{tab:layers}
\small
\begin{tabularx}{\textwidth}{@{}l X X@{}}
\toprule
\textbf{Layer} & \textbf{Responsibility} & \textbf{Principal components} \\
\midrule
Worker &
  Represent each ISA-95 Level~3 function as an autonomous BDI agent;
  translate local beliefs and domain objectives into utility-annotated
  proposals over the joint decision space. &
  Four BDI agents: \PSA{} (production), \PMA{} (maintenance),
  \QCA{} (quality), \SCA{} (supply). \\
\addlinespace
Marketplace &
  Aggregate the agents' proposals for each cross-functional conflict and
  resolve them into one composite contract that maximises joint
  welfare. &
  Bid Bus, Mediator (admissible-heuristic search), Contract Record. \\
\addlinespace
Fence &
  Audit every resolved contract against pre-declared safety rules and
  classify it into a three-tier severity scheme; emit an audit stream.
  Log-only within the scope of this thesis. &
  Rule suite, three-tier classifier, audit record stream. \\
\addlinespace
School &
  Provide the architectural slot for continual adaptation of agent
  preference weights from operational history. Deferred as future work
  (Section~\ref{sec:arch-school}). &
  Weight-adaptation module (specified, not implemented). \\
\bottomrule
\end{tabularx}
\end{table}

The layering admits a compact reading. The Worker layer is where
\emph{domain knowledge} lives: each agent knows its own function and
nothing else. The Marketplace layer is where \emph{coordination}
happens: it is informationally limited --- it sees only what the agents
publish --- but computationally active, since resolving a conflict
requires a non-trivial search. The Fence layer is where
\emph{governance} happens: it neither participates in coordination nor
holds domain knowledge; it observes the Marketplace's output and judges
it. The School layer is where \emph{adaptation} would happen: it closes
a slow outer loop from operational outcomes back to agent preferences.
Domain knowledge, coordination, governance and adaptation are thus four
separable concerns, and the architecture devotes exactly one layer to
each.

\subsection{Coordination events and the information flow}
\label{subsec:info-flow}

\NEXUS{} is event-driven. The plant floor is modelled as a discrete-event
simulation in which the four functions advance their own state until a
\emph{cross-functional conflict} arises --- a situation in which two or
more functions have a stake in the same upcoming decision. Three
conflict types are distinguished, named for the function that first
detects them: a \textsc{maintenance} conflict (a machine's modelled
health has fallen below the preventive-maintenance threshold), a
\textsc{quality} conflict (an observed defect rate has crossed a control
limit), and a \textsc{supply} conflict (a spare-part lead time has
exceeded the shortage horizon). Every conflict, whatever its type,
triggers the same coordination event, and every coordination event
produces exactly one composite contract.

A coordination event proceeds in five phases, illustrated in
Figure~\ref{fig:arch-flow}.

\begin{enumerate}[leftmargin=1.8em,itemsep=0.25em]
  \item \textbf{Announcement.} The conflict is published on the Bid Bus
    as a \emph{conflict announcement} carrying the conflict type, the
    machine identifier, a timestamp, and the small set of observable
    signals the agents are permitted to read (machine health, defect
    rate, spare lead time). The announcement carries no agent's private
    state.
  \item \textbf{Bidding.} Each of the four agents independently reads
    the announcement, updates its beliefs, derives its private
    constraint set, and emits a set of \emph{bids} ---
    utility-annotated regions of the joint contract space
    (Section~\ref{sec:arch-marketplace}).
  \item \textbf{Mediation.} The Mediator collects the four bid sets and
    searches for the combination of bids --- one per agent --- whose
    regions have a non-empty joint intersection and whose utilities sum
    to the maximum achievable value.
  \item \textbf{Contract selection.} Within the joint intersection of
    the winning bid combination the Mediator selects a single composite
    contract and publishes it back to all four agents, which then act
    on it.
  \item \textbf{Audit.} The Contract Record appends the resolved
    contract --- together with the bids that produced it --- to an
    append-only log. The Fence layer reads this log and classifies the
    contract into one of three safety tiers.
\end{enumerate}

\nexusfigure{fig:arch-flow}{Information flow of a single coordination
event. The five phases run left to right; the Fence audit in phase~5 is
post-hoc and does not feed back into the negotiation.}{a
sequence diagram with four agent lifelines (PSA, PMA, QCA, SCA), a Bid
Bus lifeline, a Mediator lifeline, a Contract Record lifeline and a
Fence lifeline. Messages, top to bottom: conflict announcement broadcast
from Bid Bus to all four agents; four bid-set messages from the agents
to the Mediator; an internal ``branch-cut A* search'' self-message on
the Mediator; a composite-contract message from the Mediator broadcast
to the four agents; an append message from the Mediator to the Contract
Record; an audit-log read from the Contract Record to the Fence; a
self-message ``three-tier classification'' on the Fence.}

Two properties of this flow are load-bearing for the rest of the
chapter. First, the negotiation is \emph{single-round}: agents bid
once, the Mediator resolves once, and the event closes. The protocol of
Ito, Hattori and Klein admits a multi-round iterative-narrowing
extension \cite{hattori2007narrowing}, in which the Mediator returns
partial-agreement feedback and agents resubmit refined bids; \NEXUS{}
deliberately implements only the single-round form and treats iterative
narrowing as future work (the justification is given in
Section~\ref{subsec:protocol-scope}). Second, the Fence audit in
phase~5 is strictly \emph{post-hoc}: it observes the contract after it
has been published and acted upon, and within the scope of this thesis
it does not block, alter, or roll back the contract. The consequences of
this log-only stance, and the path to active enforcement, are the
subject of Section~\ref{sec:arch-fence}.

\subsection{Why a layered, distributed architecture}
\label{subsec:why-layered}

Before describing the layers in detail it is worth making explicit why
the architecture is layered and distributed at all, rather than
centralised. Three arguments support the choice.

The first is \emph{organisational fidelity}. ISA-95 Level~3 is already
decomposed into four functions, each with its own information system,
its own staff and its own performance criteria \cite{sprock2019isa95}.
An architecture whose components map one-to-one onto those functions
can be deployed incrementally, function by function, and can be reasoned
about by the people who own each function. A monolithic optimiser has
no such correspondence; it must be adopted all at once and is owned by
no one.

The second is \emph{informational feasibility}. A centralised optimiser
needs the joint problem in full: every function's cost structure, every
machine's degradation model, every supplier's reliability distribution.
In a real plant these are held in separate, heterogeneous systems and
are, for organisational reasons, not freely shared. \NEXUS{} requires
only that each function be able to express its preferences as bids over
a shared, low-dimensional decision space --- a far weaker requirement
than full disclosure (Section~\ref{sec:arch-privacy}).

The third is \emph{separation of concerns}. Coordinating four functions,
governing the safety of the coordinated decision, and adapting the
functions' preferences over time are three different problems with three
different rates of change and three different owners. Folding them into
one component couples them; giving each its own layer keeps them
independently replaceable. The Fence layer, in particular, can be
strengthened from log-only auditing to active enforcement
(Section~\ref{sec:arch-fence}) without touching the Marketplace, exactly
because it is a separate layer that consumes only an audit log.

These arguments place \NEXUS{} in the established tradition of
reference architectures for manufacturing control. PROSA
\cite{vanbrussel1998prosa} introduced the idea of decomposing plant
control into cooperating holons with distinct responsibilities; ADACOR
\cite{leitao2007holonic} added an adaptive coordination mechanism that
shifts between a more centralised and a more distributed regime; HAPBA
\cite{leitao2015hapba} and the broader holonic and multi-agent
manufacturing literature \cite{leitao2009survey,trentesaux2009,monostori2006}
continued the line. \NEXUS{} inherits the layered, distributed stance of
this tradition and contributes a specific, well-characterised choice for
the coordination layer: multi-issue mediated negotiation in place of the
contract-net or hierarchical-dispatch mechanisms used by its
predecessors.

% ------------------------------------------------------------
\section{The Worker Layer}
\label{sec:arch-worker}

\subsection{Four agents, four functions}
\label{subsec:four-agents}

The Worker layer contains exactly four agents, one for each ISA-95
Level~3 operations-management function:

\begin{description}[leftmargin=1.4em,style=nextline,itemsep=0.3em]
  \item[\PSA{} --- Production Scheduling Agent.] Owns the production
    function. Its concern is throughput and the timely completion of
    operations; its home decision is the production \emph{rate} at which
    a machine is driven.
  \item[\PMA{} --- Preventive Maintenance Agent.] Owns the maintenance
    function. Its concern is machine health and the avoidance of
    unplanned breakdowns; its home decision is the maintenance
    \emph{window} --- whether and how urgently a machine is serviced.
  \item[\QCA{} --- Quality Control Agent.] Owns the quality function.
    Its concern is the defect rate and the cost of defect-driven
    rework; its home decision is the quality \emph{action} --- the
    intensity of any inspection hold placed on output.
  \item[\SCA{} --- Supply Chain Agent.] Owns the inventory and supply
    function. Its concern is spare-part availability and stockout
    exposure; its home decision is the supply \emph{horizon} --- how far
    ahead spare parts are preordered.
\end{description}

The cardinality is fixed by the ISA-95 decomposition itself, not chosen
for convenience: Level~3 defines four operations-management functions,
so \NEXUS{} has four Worker agents. Table~\ref{tab:agents} records the
mapping and previews the home issue each agent decides --- the issues
themselves are defined in Section~\ref{sec:arch-marketplace}.

\begin{table}[ht]
\centering
\caption{The four Worker agents, their ISA-95 Level~3 functions, their
local beliefs, and the home issue each one owns in the composite
contract.}
\label{tab:agents}
\small
\begin{tabularx}{\textwidth}{@{}l l X l@{}}
\toprule
\textbf{Agent} & \textbf{ISA-95 L3 function} & \textbf{Principal local beliefs}
  & \textbf{Home issue} \\
\midrule
\PSA{} & Production operations & Operation backlog, machine load, tardiness
  pressure & \texttt{production\_rate} \\
\PMA{} & Maintenance operations & Machine health, time since last service,
  breakdown hazard & \texttt{maint\_window} \\
\QCA{} & Quality operations & Observed defect rate, severity, rework backlog
  & \texttt{quality\_action} \\
\SCA{} & Inventory / supply operations & Spare-part stock, supplier lead time
  & \texttt{supply\_horizon} \\
\bottomrule
\end{tabularx}
\end{table}

\subsection{The BDI agent model}
\label{subsec:bdi}

Each Worker agent is structured according to the
Belief--Desire--Intention model of agency \cite{rao1995bdi,wooldridge1995}.
The BDI model is a natural fit for an ISA-95 Level~3 function because
the three mental attitudes map cleanly onto the three things a function
must do: maintain a picture of its part of the plant, hold objectives,
and commit to actions.

\begin{description}[leftmargin=1.4em,style=nextline,itemsep=0.3em]
  \item[Beliefs.] The agent's local model of its slice of the plant.
    \PSA{} believes things about operation backlogs and machine load;
    \PMA{} about machine health and time since service; \QCA{} about
    defect rates and severity; \SCA{} about stock levels and supplier
    lead times. Beliefs are updated from the observations the agent is
    entitled to read and from the conflict announcement. Crucially,
    \emph{an agent's beliefs are private}: no agent observes another
    agent's beliefs, and the Marketplace never reads them
    (Section~\ref{sec:arch-privacy}).
  \item[Desires.] The agent's objectives, expressed as a small set of
    weighted goals within its own function --- for \PSA{}, for example,
    minimising tardiness, maximising throughput, and minimising
    disruption. The desire weights are parameters of the agent. Within
    the scope of this thesis they are fixed at calibrated values; the
    School layer (Section~\ref{sec:arch-school}) is the architectural
    slot in which they would later be adapted.
  \item[Intentions.] The agent's committed proposals for the current
    coordination event. In \NEXUS{} an agent's intentions materialise
    as a set of \emph{bids} --- the utility-annotated regions of the
    joint contract space defined in Section~\ref{sec:arch-marketplace}.
    The bid set is the only thing an agent makes public; it is the
    externally visible projection of the agent's beliefs and desires.
\end{description}

Figure~\ref{fig:arch-bdi} shows the internal structure of a Worker
agent and the deliberation pipeline that turns a conflict announcement
into a bid set. The pipeline has two stages the agent specialises and
one stage shared by all four agents. The specialised stages are belief
update --- how the agent revises its local model from a new observation
--- and constraint derivation --- how the agent's current beliefs and
desires generate a private set of weighted constraints over the joint
contract space. The shared stage is bid generation, which runs the same
four-step procedure for every agent (Section~\ref{subsec:bid-gen}).

\nexusfigure{fig:arch-bdi}{Internal structure of a Worker agent. The
belief-update and constraint-derivation stages are specialised per
agent; the bid-generation stage is shared. Only the bid set crosses the
agent boundary.}{a single agent box, internally divided into
``Beliefs'' (a small database icon), ``Desires'' (a weighted goal list),
and a deliberation pipeline with three stages drawn left to right:
``update beliefs'' (fed by an incoming observation arrow),
``derive constraints'' (producing a stack of weighted constraint icons),
and ``generate bids'' (producing bid icons). An incoming arrow on the
left is labelled ``conflict announcement''; the only outgoing arrow, on
the right, is labelled ``bid set'' and crosses the dashed agent
boundary. Beliefs and Desires are shown entirely inside the boundary,
with a small lock icon indicating they never cross it.}

\subsection{Mono-dimensional outputs and domain authority}
\label{subsec:mono-dim}

A central discipline of the Worker layer is that each agent's authority
is \emph{mono-dimensional}: an agent is the owner of exactly one issue
of the composite contract --- its home issue --- and although it is
permitted to express preferences over the other three issues, those
cross-domain preferences are constructed to be, by design, an order of
magnitude weaker than its preferences on its home issue.

This discipline answers a question that the multi-issue setting raises
immediately. If every agent may bid on every issue, what stops \QCA{}
from effectively dictating the maintenance decision, or \PSA{} from
overriding a quality hold? In a single-winner mechanism the answer
would have to be a separate priority rule layered on top of the
mechanism. In \NEXUS{} the answer is structural and lives inside the
bid weights: each agent's constraints on its \emph{home} issue carry
weights bounded by an agent-specific ceiling, while its constraints on
\emph{other} issues carry weights produced by multiplying a base weight
by a small cross-domain damping factor. The home-issue ceilings ---
recorded in Table~\ref{tab:ceilings} --- are roughly an order of
magnitude above the typical damped cross-domain weight. The encoding
realises ISA-95 functional authority operationally: a function's
preference on its own decision dominates the cross-domain stakes of its
peers, so the maintenance decision is, in the normal case, \PMA{}'s to
make --- but the peers still have a voice, and that voice shapes the
outcome whenever \PMA{} itself is indifferent between maintenance
options.

\begin{table}[ht]
\centering
\caption{Approximate per-agent home-issue weight ceilings (domain
authority). Cross-domain constraint weights are damped to roughly a
tenth of these values, so an agent's preference on its home issue
dominates its peers' cross-domain stakes by about an order of
magnitude.}
\label{tab:ceilings}
\small
\begin{tabular}{@{}l l r@{}}
\toprule
\textbf{Agent} & \textbf{Home issue} & \textbf{Home-issue weight ceiling} \\
\midrule
\PMA{} & \texttt{maint\_window}     & $\approx 500$ \\
\PSA{} & \texttt{production\_rate}  & $\approx 200$ \\
\SCA{} & \texttt{supply\_horizon}   & $\approx 200$ \\
\QCA{} & \texttt{quality\_action}   & $\approx 175$ \\
\bottomrule
\end{tabular}
\end{table}

It is worth being precise about what ``mono-dimensional'' does and does
not mean. It does \emph{not} mean an agent estimates only its own
issue: each agent emits constraints over the whole joint space, and the
cross-domain constraints are what make the negotiation genuinely
multi-issue rather than four parallel single-issue negotiations. It
\emph{does} mean three things. First, no agent observes another agent's
domain state --- \PSA{} never sees \PMA{}'s health model, only the
shared signals carried in the announcement; there is no inter-domain
estimation and no cross-visibility of beliefs. Second, an agent's
cross-domain preferences are subordinate by construction, through the
damping factor, so they act as tie-breakers and nudges rather than
overrides. Third, ownership is exclusive: each issue has exactly one
owning agent, and the architecture never asks two agents to arbitrate
which of them owns a decision. The first point is an informational
property and underpins the privacy claim of
Section~\ref{sec:arch-privacy}; the second and third are authority
properties and underpin the claim that \NEXUS{} respects the ISA-95
functional decomposition rather than dissolving it.

The empirical behaviour of the weight ceilings is discussed in
Chapter~\ref{ch:results}; here it is enough to note that on the
benchmark families studied the order-of-magnitude separation makes the
exact cross-domain damping factor a low-sensitivity parameter --- the
qualitative outcome is stable across a range of damping values --- which
is the expected signature of a design in which cross-domain preferences
are deliberately subdominant.

% ------------------------------------------------------------
\section{The Marketplace Layer}
\label{sec:arch-marketplace}

The Marketplace layer is where coordination happens. It receives the
four agents' bid sets for a conflict and returns one composite contract.
This section defines the composite contract space
(Section~\ref{subsec:contract-space}), the bid representation and the
agent-side procedure that produces it (Sections~\ref{subsec:bids}
and~\ref{subsec:bid-gen}), the mediator's search for the
welfare-maximising joint contract (Section~\ref{subsec:mediator}), and
the scope decisions that bound the protocol (Section~\ref{subsec:protocol-scope}).

\subsection{The composite contract space}
\label{subsec:contract-space}

The defining decision of the Marketplace layer --- and the decision
that most sharply distinguishes \NEXUS{} from a single-winner
coordination mechanism --- is that a coordination event is resolved not
by selecting one agent's preferred action but by producing a
\emph{composite contract}: a vector that assigns a value to every one of
the four functional decisions simultaneously.

Formally, the composite contract space is the Cartesian product of four
\emph{issue} domains, one per agent home issue:
%
\begin{equation}
  \label{eq:contract-space}
  \Omega \;=\; \Omega_{\mathrm{PR}} \times \Omega_{\mathrm{MW}}
            \times \Omega_{\mathrm{QA}} \times \Omega_{\mathrm{SH}},
\end{equation}
%
where $\Omega_{\mathrm{PR}}$ is the domain of the production rate,
$\Omega_{\mathrm{MW}}$ of the maintenance window, $\Omega_{\mathrm{QA}}$
of the quality action, and $\Omega_{\mathrm{SH}}$ of the supply horizon.
A \emph{contract} is a single point $c \in \Omega$. The four issue
domains, with their discrete levels, are given in
Table~\ref{tab:issues}.

\begin{table}[ht]
\centering
\caption{The four issues of the composite contract and their discrete
domains. The joint space has $|\Omega| = 3 \times 3 \times 3 \times 4 =
108$ contracts.}
\label{tab:issues}
\small
\begin{tabularx}{\textwidth}{@{}l l l X@{}}
\toprule
\textbf{Issue} & \textbf{Owner} & \textbf{Levels} & \textbf{Interpretation} \\
\midrule
\texttt{production\_rate} & \PSA{} & \textsc{slow}, \textsc{normal}, \textsc{fast}
  & Speed at which the machine is driven; faster lifts output but
    accelerates wear. \\
\addlinespace
\texttt{maint\_window} & \PMA{} & \textsc{none}, \textsc{planned}, \textsc{immediate}
  & Whether and how urgently the machine is serviced. \\
\addlinespace
\texttt{quality\_action} & \QCA{} & \textsc{no\_hold}, \textsc{hold\_light}, \textsc{hold\_full}
  & Intensity of the inspection hold placed on output. \\
\addlinespace
\texttt{supply\_horizon} & \SCA{} & \textsc{h0}, \textsc{h24}, \textsc{h48}, \textsc{h72}
  & Length, in hours, of the spare-part preorder window. \\
\bottomrule
\end{tabularx}
\end{table}

The joint space is small --- $108$ contracts --- and finite. This is a
deliberate scope decision. The issue domains are kept discrete and
coarse so that the joint space can be enumerated exactly and the
mediator's search can carry a global-optimality guarantee
(Section~\ref{subsec:mediator}); the alternative of continuous or
finely discretised issues would force the mediator onto approximate
search and would forfeit that guarantee for no benefit at the scale of
an ISA-95 Level~3 coordination loop.

The composite structure is what makes the cross-functional trade-offs
representable. A single-winner mechanism resolves a conflict by choosing
a label --- ``maintenance wins'' --- and so can never express the
statement ``service the machine \emph{and} hold quality light
\emph{and} keep the rate normal''; the composite contract expresses
exactly such joint statements, because it is a point in
$\Omega$ rather than a choice among the agents. It is structurally
impossible for the mechanism to collapse into fixed-priority routing,
because its output is a vector by construction. This point is developed
further when the contrast with single-winner mechanisms is drawn in
Section~\ref{subsec:protocol-scope}.

\subsection{Bids: utility-annotated regions}
\label{subsec:bids}

Agents do not reveal their utility functions to the Marketplace. They
reveal \emph{bids}. Following the multi-issue protocol of Ito, Hattori
and Klein \cite{ito2007multiissue}, a bid is a compact, public summary
of one region of the contract space that the agent finds acceptable,
annotated with the utility the agent assigns to it.

An agent $i$ holds, for the current conflict, a private set of weighted
\emph{constraints}. A constraint $k$ is a predicate over the contract
space --- it is satisfied by a subset $R_{ik} \subseteq \Omega$ of
contracts --- carrying a non-negative weight $w_{ik}$. The agent's
utility for a contract $c$ is the additive sum of the weights of the
constraints that $c$ satisfies,
%
\begin{equation}
  \label{eq:utility}
  U_i(c) \;=\; \sum_{k \,:\, c \,\in\, R_{ik}} w_{ik},
\end{equation}
%
which is the additive utility form of Ito, Hattori and Klein
\cite{ito2007multiissue} and is the natural multi-issue generalisation
of the weighted-goal scoring used in single-issue agent bidding.
A \emph{bid} $b_{ij}$ submitted by agent $i$ is then the pair
%
\begin{equation}
  \label{eq:bid}
  b_{ij} \;=\; \bigl(\, u_{ij},\; R_{ij} \,\bigr),
\end{equation}
%
where $R_{ij} \subseteq \Omega$ is a region of the contract space ---
in practice a hyper-rectangle, the intersection of the regions of the
agent's most heavily weighted satisfied constraints --- and $u_{ij}$ is
the scalar utility the agent attaches to that region. The agent submits
a small set of such bids per conflict.

The width of the bid region is governed by a relaxation parameter,
written \textsc{top}-\textit{n}: the region $R_{ij}$ is formed as the
intersection of the $n$ most heavily weighted constraints the agent
satisfies, while the annotated utility $u_{ij}$ aggregates the weights
of \emph{all} satisfied constraints. The parameter trades expressiveness
against feasibility. With $n=1$ the region is just the single heaviest
constraint's region --- wide, easy to intersect with other agents'
regions, but blind to the agent's inter-issue trade-offs. With larger
$n$ the region is the intersection of more constraints --- narrower,
more faithful to the agent's joint preference, but more likely to leave
no contract in common with the other agents, which would cause the
negotiation to fail. \NEXUS{} uses $n=2$: the bid region is the
intersection of an agent's two heaviest constraints, which captures the
single most important inter-issue trade-off each agent holds while
keeping the empirical negotiation-failure rate low. The calibration
behind this choice, and the failure-rate-versus-expressiveness curve it
sits on, are reported in Chapter~\ref{ch:results}.

What an agent publishes is therefore exactly its bid set --- a bounded
number of (utility, region) pairs --- and nothing else. It does not
publish the constraints themselves, their weights, the beliefs that
generated them, or the desire weights that scaled them. This is the
precise sense in which \NEXUS{} preserves informational locality, and
Section~\ref{sec:arch-privacy} states the property formally.

\subsection{Bid generation}
\label{subsec:bid-gen}

Each agent turns its private constraint set into a bid set by a
four-step procedure, shared by all four agents and adapted from the
sampling-and-adjustment scheme of the multi-issue protocol
\cite{ito2007multiissue,klein2003negotiating}:

\begin{enumerate}[leftmargin=1.8em,itemsep=0.2em]
  \item \textbf{Sample.} Draw a set of candidate contracts from the
    joint space $\Omega$.
  \item \textbf{Adjust.} From each sampled contract, perform a local
    search --- in \NEXUS{} a greedy hill-climb --- towards a nearby
    contract of higher private utility $U_i$.
  \item \textbf{Bid.} For each adjusted contract whose utility exceeds
    a reservation threshold, build a bid: the region is the
    \textsc{top}-\textit{n} constraint intersection around the contract, the
    utility is the value of Equation~\eqref{eq:utility}.
  \item \textbf{Submit.} Emit the resulting set of bids to the Bid Bus.
\end{enumerate}

The local-search step admits two implementations. The original protocol
uses simulated annealing, accepting an uphill or a controlled downhill
move with the Metropolis probability
$\exp\bigl((U_i(c') - U_i(c))/T\bigr)$ at temperature $T$
\cite{klein2003negotiating}. \NEXUS{} retains simulated annealing as a
selectable option but defaults to plain greedy hill-climbing. The reason
is empirical and is documented in
Chapter~\ref{ch:implementation}: at the scale of the \NEXUS{} contract
space the annealing temperature has no measurable effect on the
resulting bids, so the simpler, faster, fully deterministic greedy
search is preferred. Determinism is itself valuable here --- it is part
of what makes a coordination event byte-for-byte reproducible, which the
reproducibility infrastructure of Chapter~\ref{ch:implementation} relies
on.

Listing~\ref{lst:bid} sketches the data structures involved, in the
notation of the implementation.

\begin{lstlisting}[language=Python,caption={The constraint and bid data
structures (schematic). An agent holds a list of \texttt{Constraint}s
and emits a list of \texttt{Bid}s; only the latter is public.},label={lst:bid}]
class Constraint:        # private to the owning agent
    region: FrozenSet[Contract]   # contracts that satisfy the predicate
    weight: float                 # non-negative; scaled by desire weight

class Bid:               # public; the only thing that leaves an agent
    utility: float                # sum of ALL satisfied constraint weights
    region:  FrozenSet[Contract]  # intersection of the top-n constraints
    sample:  Contract             # a representative contract in the region
\end{lstlisting}

\subsection{The mediator}
\label{subsec:mediator}

The Mediator is the active component of the Marketplace layer. It
receives one bid set per agent and must return one composite contract.
Per the placement decision discussed in
Section~\ref{subsec:mediator-placement}, it is a component
\emph{within} the Marketplace layer, not an agent and not a supervisor
above the agents.

\paragraph{The optimisation problem.}
Let agent $i$ submit bid set $B_i$. A \emph{combination} is a selection
of one bid per agent, $\beta = (b_1, \dots, b_4)$ with $b_i \in B_i$.
The combination is \emph{feasible} if the four bid regions have a
non-empty joint intersection,
%
\begin{equation}
  \label{eq:feasible}
  R(\beta) \;=\; \bigcap_{i=1}^{4} R_{i}(b_i) \;\neq\; \emptyset,
\end{equation}
%
which is the condition that there exists at least one contract
acceptable to all four agents under the chosen bids. The Mediator
searches for the feasible combination of maximum summed utility,
%
\begin{equation}
  \label{eq:welfare}
  \beta^{\star} \;=\; \arg\max_{\beta \,:\, R(\beta) \neq \emptyset}
                      \;\sum_{i=1}^{4} u_i(b_i),
\end{equation}
%
and then commits to a single contract $c^{\star} \in R(\beta^{\star})$.
Objective~\eqref{eq:welfare} is a social-welfare maximisation: it picks
the joint agreement that maximises the sum of the agents' utilities.
Under the cooperative, additive-utility setting of
Equation~\eqref{eq:utility}, the maximiser of the sum cannot be
Pareto-dominated --- any contract that improved one agent without
harming another would have a strictly greater sum and so would already
have been chosen. The composite contract \NEXUS{} produces is therefore
\emph{Pareto-optimal with respect to the published bids} by
construction. This property is load-bearing for the comparison against
a centralised Pareto optimiser in Chapter~\ref{ch:results}.

\paragraph{Why search rather than enumerate.}
The space of combinations is bounded --- with four agents and at most
$K$ bids each there are at most $K^4$ of them --- and at the
calibrated bid counts of \NEXUS{} a brute-force enumeration would in
fact be tractable. The Mediator nonetheless uses a search, for two
reasons. The first is efficiency: most combinations are infeasible, and
their infeasibility can be detected on a \emph{partial} selection ---
once two bid regions fail to intersect, no extension of that partial
selection can intersect either. A search that builds combinations one
agent at a time can therefore prune (``branch-cut'') an infeasible
partial selection before ever extending it, and so does work
proportional to the feasible part of the search tree rather than to
$K^4$. The second reason is that, with the right priority ordering, the
same search yields a provable global-optimality guarantee at its first
complete result --- which a fixed enumeration order does not.

\paragraph{Branch-cut search with an admissible heuristic.}
The Mediator performs a best-first search over partial combinations,
expanding one agent at a time. Each partial combination $p$ at depth
$d$ (i.e.\ with $d$ of the four agents' bids chosen) is placed on a
priority queue keyed by
%
\begin{equation}
  \label{eq:astar}
  f(p) \;=\; \underbrace{\sum_{i \le d} u_i(b_i)}_{g(p)\ :\ \text{utility committed}}
           \;+\; \underbrace{\sum_{j > d} \max_{b \in B_j} u(b)}_{h(p)\ :\ \text{optimistic remainder}} ,
\end{equation}
%
and partial combinations whose joint region has become empty are
discarded (branch-cut) rather than queued. The term $g(p)$ is the
utility already committed by the chosen bids; the term $h(p)$ is, for
each not-yet-chosen agent, the largest utility any of its bids could
contribute. Because branch-cutting can only ever \emph{remove}
feasible completions, $h(p)$ never underestimates the true utility
reachable from $p$: it is an \emph{admissible} heuristic in the sense of
the A* algorithm of Hart, Nilsson and Raphael. Expanding the queue in
decreasing order of $f$ is therefore A* search on a maximisation
problem, and the standard A* optimality argument applies: \emph{the
first complete combination the search returns is a global maximiser of
Equation~\eqref{eq:welfare}}.

The guarantee is not a cosmetic refinement. An earlier version of the
Mediator keyed the queue on the committed utility $g(p)$ alone, omitting
the optimistic remainder $h(p)$. That weaker key makes the search
best-first on partial sums but \emph{not} globally optimal: it can
commit to a high-utility first bid all of whose feasible completions are
poor, and never explore a lower-utility first bid with an excellent
forced completion. On adversarially constructed inputs the resulting
contract was found to be up to $85\%$ short of the optimum. Adding the
admissible term $h(p)$ closes that gap and aligns the implementation
with the social-welfare claim the architecture makes. The episode is
recorded here because it illustrates a general point about reference
architectures: a claimed property --- here, ``the mediator maximises
joint welfare'' --- is only as good as the algorithm that is supposed
to deliver it, and the two must be verified to agree.

Figure~\ref{fig:arch-mediator} illustrates the search on a small
example. Algorithm~\ref{lst:mediator} gives the procedure in pseudocode.

\nexusfigure{fig:arch-mediator}{Branch-cut A* search over partial bid
combinations. Each level of the tree fixes one more agent's bid; a
crossed node is a branch-cut (empty joint region); the highlighted path
is the first complete combination popped, which is the global welfare
maximiser.}{a four-level search tree. The root is labelled
``no bids chosen''. Level~1 nodes fix a \PMA{} bid, level~2 a \QCA{}
bid, level~3 a \PSA{} bid, level~4 an \SCA{} bid. Several interior nodes
are marked with a cross and the label ``branch-cut: empty joint
region''. Each node carries a small label ``$f = g + h$''. One
root-to-leaf path is drawn in bold and labelled ``first complete pop ---
global welfare maximiser''. A side caption notes that bold-path leaf
utility equals $\max \sum u_i$.}

\begin{lstlisting}[caption={The mediator's branch-cut A* search
(pseudocode). The priority key f = g + h is Equation~(3.6).},label={lst:mediator}]
function SELECT_CONTRACT(bid sets B_1 .. B_4):
    queue <- empty priority queue, ordered by f, largest first
    push the empty partial combination, with f = sum_j ( max utility in B_j )
    while queue is not empty:
        p <- pop the partial combination of largest f
        if p is complete (one bid chosen for all 4 agents):
            return a contract c* from the joint region R(p)   # global optimum
        d <- number of bids already chosen in p
        for each bid b in B_(d+1):
            p' <- p extended with bid b
            if the joint region R(p') is empty:
                discard p'                                    # branch-cut
            else:
                push p' with key f(p') = g(p') + h(p')        # see Eq. 3.6
    return FAILURE   # no feasible combination exists
\end{lstlisting}

\paragraph{Contract selection within the joint region.}
The search returns a winning combination $\beta^{\star}$; its joint
region $R(\beta^{\star})$ typically still contains more than one
contract. The Mediator selects the final $c^{\star}$ from that region by
a deterministic tie-break --- maximising the summed agent utility over
the contracts in the region, and, where that is itself tied, taking the
lowest-indexed contract under a fixed enumeration order. Determinism in
the tie-break is what makes the whole coordination event reproducible:
the same bid sets always yield the same contract.

\paragraph{Negotiation failure.}
If every combination is infeasible --- the search empties its queue
without ever popping a complete combination --- the event is a
\emph{negotiation failure}: no contract is acceptable to all four
agents under their submitted bids. The environment then falls back to a
conservative default, and, as Section~\ref{sec:arch-fence} describes,
the Fence layer records the failure as a top-tier audit signal. The
failure rate is a calibration concern --- it is the quantity the
\textsc{top}-\textit{n} parameter trades against expressiveness --- and is
reported in Chapter~\ref{ch:results}.

\subsection{Mediator placement and protocol scope}
\label{subsec:protocol-scope}

\paragraph{The mediator is inside the Marketplace, not above the agents.}
\label{subsec:mediator-placement}
It would be possible to model the Mediator as a fifth agent --- a
``hyper-agent'' sitting above \PSA{}, \PMA{}, \QCA{} and \SCA{} and
arbitrating between them. \NEXUS{} deliberately does not. The Mediator
is a logical component \emph{within} the Marketplace layer for a precise
reason: it has no beliefs, no desires and no model of the plant. It
does not know what a machine is. Its entire information set is the union
of the bids published for the current conflict, and its entire function
is the combinatorial search of Equation~\eqref{eq:welfare} over that
information set. A fifth BDI agent, by contrast, would have its own
utility function and its own view of the plant state, and would
re-introduce exactly the centralised problem-knowledge that the
distributed architecture exists to avoid. Modelling the Mediator as an
in-layer component keeps the architectural diagram honest: \NEXUS{} has
four agents, not five, and the coordination layer is a mechanism, not a
master.

\paragraph{Single-round protocol.}
\NEXUS{} implements the single-round form of the multi-issue protocol:
agents bid once, the Mediator resolves once. The iterative-narrowing
extension \cite{hattori2007narrowing}, in which the Mediator returns
partial-agreement feedback and agents resubmit refined bids over
several rounds, is not implemented. The decision is a scope decision
and is justified by the scale of the problem. Iterative narrowing earns
its cost on problems where the number of agents or the per-agent bid
count is large enough that the single-round combination space becomes
intractable; \NEXUS{} has exactly four agents --- fixed by the ISA-95
decomposition --- and a small per-agent bid count, so the single-round
search is comfortably tractable and carries the global-optimality
guarantee of Section~\ref{subsec:mediator}. Implementing iterative
narrowing would add round management, convergence criteria and an
agent-side refinement strategy for no gain at the demonstrated scale.
The extension is a clean future-work item and is revisited in
Chapter~\ref{ch:conclusions}.

\paragraph{Why multi-issue mediation rather than a single-winner mechanism.}
The Marketplace mechanism was selected from a field of candidate
multi-agent coordination paradigms; Table~\ref{tab:paradigms}
summarises the comparison, which Chapter~\ref{ch:literature} develops in
full. The decisive consideration is the one already met in
Section~\ref{subsec:contract-space}. A single-winner mechanism such as
the Contract Net Protocol \cite{smith1980,deen1997cnp} resolves a
conflict by awarding it to one bidder; its output is a label, and a
label cannot express a cross-functional trade-off. Distributed
constraint optimisation \cite{modi2005adopt,fioretto2018dcop} does
produce a joint assignment and does so with optimality guarantees, but
it requires agents to publish full cost tables over the shared
variables --- precisely the disclosure that the bounded-disclosure
commitment of Section~\ref{sec:arch-intro} forbids. Bilateral
negotiation \cite{faratin1998negotiation} scales poorly to four parties
and has no native representation of joint trade-offs across more than
two issues; generalised partial global planning \cite{decker1995gpgp}
and argumentation-based negotiation \cite{rahwan2003argumentation} were
set aside on grounds of implementation complexity and runtime cost
respectively. Multi-issue multi-lateral mediation
\cite{ito2007multiissue,fatima2004multiissue} is the design point that
satisfies all of \NEXUS{}'s requirements at once: its output is a
composite contract, its disclosure is bounded to bid summaries, and it
is tractable at four agents.

\begin{table}[ht]
\centering
\caption{Coordination paradigms considered for the Marketplace layer.
Only multi-issue mediation produces a composite contract, bounds
disclosure to bid summaries, and is tractable at four agents.
Chapter~\ref{ch:literature} develops the comparison.}
\label{tab:paradigms}
\small
\begin{tabularx}{\textwidth}{@{}l c c X@{}}
\toprule
\textbf{Paradigm} & \textbf{Composite} & \textbf{Bounded} & \textbf{Status in \NEXUS{}} \\
 & \textbf{output} & \textbf{disclosure} & \\
\midrule
Contract Net (single-winner) & no & yes &
  Rejected: label output cannot express trade-offs. \\
Bilateral negotiation & partial & yes &
  Rejected: poor scaling to four parties. \\
GPGP & yes & partial &
  Rejected: implementation complexity. \\
DCOP & yes & \textbf{no} &
  Rejected: requires full cost-table disclosure. \\
Argumentation-based & yes & yes &
  Rejected: runtime cost, ontology requirement. \\
\textbf{Multi-issue mediation} & \textbf{yes} & \textbf{yes} &
  \textbf{Selected} (Ito, Hattori \& Klein). \\
\bottomrule
\end{tabularx}
\end{table}

% ------------------------------------------------------------
\section{The Fence Layer}
\label{sec:arch-fence}

\subsection{Governance as a separate concern}
\label{subsec:fence-rationale}

The Marketplace layer maximises joint welfare over the agents'
published bids. It does not, and structurally cannot, guarantee that the
resulting contract is \emph{safe}. Welfare maximisation is an
optimisation over preferences; safety is a constraint on outcomes, and
the two are different things. A contract that runs a machine at the
fast rate when its modelled health is already near the breakdown
threshold may well be the joint-welfare maximiser for a given set of
bids --- if, for instance, \PMA{}'s constraints happened not to fire
strongly on that conflict --- yet it is plainly unsafe. Some component
must judge the contract against safety criteria that are independent of
the agents' preferences.

The Fence layer is that component. It is deliberately a \emph{separate
layer} rather than a check folded into the Mediator, for the
separation-of-concerns reason given in
Section~\ref{subsec:why-layered}: safety rules change at a different
rate, and are owned by different people, than the coordination
mechanism. Folding the rules into the Mediator would couple the two and
would make the Mediator's welfare guarantee conditional on the current
rule set. Keeping the Fence separate, and having it consume only the
Marketplace's audit log, keeps the Mediator's guarantee clean and lets
the safety rules evolve independently. This structural integration of
safety governance --- as an audit layer over the coordination
substrate, with no protocol-level coupling --- is one of the
architectural claims of the thesis.

\subsection{Three-tier classification}
\label{subsec:fence-tiers}

The Fence holds a suite of \emph{pre-declared safety rules}. Each rule
is a pair of predicates over a contract and the plant signals
associated with its conflict: a \emph{violation} predicate, true when
the contract hard-violates the rule, and a \emph{boundary} predicate,
true when the contract does not yet violate the rule but lies within a
narrow band --- ten per cent of the safe range --- of doing so. For
every contract in the Contract Record the Fence evaluates every rule
and assigns a severity tier:

\begin{description}[leftmargin=2.6em,style=nextline,itemsep=0.25em]
  \item[$\sigma_0$ --- Normal.] No rule is violated and no rule is at
    its boundary. The contract is safe.
  \item[$\sigma_1$ --- Warning.] No rule is violated, but at least one
    rule is within its boundary band. The contract is safe but is
    approaching a safety limit.
  \item[$\sigma_2$ --- Critical.] At least one rule is hard-violated.
    The contract breaches a safety limit. A negotiation failure ---
    where the Marketplace produced no contract at all --- is also
    classified $\sigma_2$ by convention: the absence of a feasible
    agreement is itself a critical audit signal.
\end{description}

The tiers are illustrated in Figure~\ref{fig:arch-fence}. The
three-tier scheme is a coarse but operationally meaningful summary: it
distinguishes contracts that are fine ($\sigma_0$), contracts that
warrant attention ($\sigma_1$), and contracts that demand it
($\sigma_2$), and it does so with a granularity a plant supervisor can
act on without further interpretation.

\nexusfigure{fig:arch-fence}{The three-tier Fence classification. Each
resolved contract is checked against every safety rule and placed in the
highest tier any rule triggers.}{a decision-flow diagram. A
contract enters from the left into a box ``evaluate all safety rules''.
Three labelled outcome bands leave the box: a green band
``$\sigma_0$ Normal --- no violation, no boundary''; an amber band
``$\sigma_1$ Warning --- a rule within 10\% of its limit''; a red band
``$\sigma_2$ Critical --- a rule hard-violated, or negotiation
failure''. A note beneath reads ``log-only: the tier is recorded, the
contract is not altered''.}

The default rule suite contains four rules, one anchored on each issue;
they are listed in Table~\ref{tab:fence-rules}. Each rule encodes a
domain safety invariant: do not run fast into an imminent breakdown, do
not skip maintenance on a critically degraded machine, do not waive a
quality hold on a severely defective batch, do not leave the supply
horizon empty when the spare-part lead time is long. The rule suite is
a parameter of the Fence: it can be extended or retuned without any
change to the other three layers, which is the practical payoff of
keeping governance in its own layer.

\begin{table}[ht]
\centering
\caption{The default Fence safety rule suite. Each rule pairs a hard
\emph{violation} predicate ($\sigma_2$) with a \emph{boundary}
predicate ($\sigma_1$) that fires within ten per cent of the safe
limit.}
\label{tab:fence-rules}
\small
\begin{tabularx}{\textwidth}{@{}l X@{}}
\toprule
\textbf{Rule} & \textbf{Violation condition ($\sigma_2$)} \\
\midrule
\texttt{no\_fast\_when\_break\-down\_imminent} &
  \texttt{production\_rate} is \textsc{fast} while machine health is
  below the breakdown threshold. \\
\addlinespace
\texttt{maintain\_when\_health\_critical} &
  \texttt{maint\_window} is \textsc{none} while machine health is below
  the preventive-maintenance threshold. \\
\addlinespace
\texttt{hold\_quality\_when\_severe} &
  \texttt{quality\_action} is \textsc{no\_hold} while the defect rate is
  at or above twice the quality control limit. \\
\addlinespace
\texttt{preorder\_when\_long\_lead} &
  \texttt{supply\_horizon} is \textsc{h0} while the spare-part lead time
  exceeds the shortage horizon. \\
\bottomrule
\end{tabularx}
\end{table}

\subsection{Log-only enforcement, and the path to active governance}
\label{subsec:fence-logonly}

Within the scope of this thesis the Fence is \emph{log-only}. It
classifies every contract and emits an audit-record stream, but it does
not block, alter, or roll back any contract --- a $\sigma_2$ contract is
executed exactly as a $\sigma_0$ one would be, and the tier is simply
recorded. The audit stream is consumed downstream: it supports the
safety analysis of Chapter~\ref{ch:results} and the deployment
discussion of Chapter~\ref{ch:industrial}.

The log-only stance is a deliberate scope boundary, and it is the right
one for two reasons. First, it lets the thesis make and defend a clean,
limited claim --- that safety governance can be \emph{integrated
structurally}, as an audit layer over the coordination substrate,
without coupling into the negotiation protocol --- and demonstrate that
claim without entangling it in the much harder questions that active
enforcement raises. Second, log-only operation is in fact how a
governance layer should be commissioned in practice: a plant would run
the Fence in observe-only mode first, accumulate evidence on how often
and why each tier fires, tune the rule suite against that evidence, and
only then enable active control. The log-only Fence of this thesis is
the first, necessary stage of exactly that commissioning path.

Active runtime enforcement --- a $\sigma_2$ contract being rejected and
sent back for re-negotiation, or overridden by a safe default --- is the
natural next step and is treated as future work in
Chapter~\ref{ch:conclusions}. It raises questions the log-only layer is
designed to sidestep: how to re-negotiate without looping, how to bound
the latency an override adds to the coordination event, and how to keep
an overriding Fence from silently becoming the de-facto coordinator. The
runtime-verification and shielding literature
\cite{raju2020runtime,alshiekh2017,konighofer2019shields} offers
mechanisms for the active case; integrating one of them into the Fence
layer is a well-defined extension, and the layer boundary --- the Fence
already consumes the contract stream and already classifies it --- is
exactly the seam at which an enforcement mechanism would attach.

% ------------------------------------------------------------
\section{The School Layer}
\label{sec:arch-school}

The fourth layer of \NEXUS{} is the School. Unlike the other three it is
\emph{not implemented} within the scope of this thesis: it is described
here, and carried in the architectural diagram with a dashed boundary,
so that the architecture is presented complete and so that the
future-work proposal of Chapter~\ref{ch:conclusions} has a concrete
home. This section states what the School layer is for and how it would
attach to the rest of the architecture; it does not claim an
implemented mechanism.

\paragraph{The problem the School layer addresses.}
The agents' desire weights --- the weights that scale each agent's
constraints, and through them its bids --- are fixed parameters in the
implemented architecture. They are calibrated once and held constant.
But a real plant is not stationary: product mixes change, machines age,
suppliers' reliability drifts, and a set of desire weights that
coordinates well in one operating regime may coordinate poorly in
another. The School layer is the architectural slot for an outer loop
that observes operational outcomes and slowly adapts the agents' desire
weights so that the coordination stays well-tuned as the plant
evolves.

\paragraph{How it would attach.}
The School layer would close a slow loop around the other three. It
would consume the same audit stream the Fence consumes --- the Contract
Record together with the realised plant outcomes --- and would adjust
each agent's desire weights in the direction that improves the observed
outcomes. The adjusted weights would be written back to the agents, and
the Marketplace would simply see better-calibrated bidders on
subsequent conflicts; no protocol change is needed, because the desire
weights are internal to the agents and the bid interface is unchanged.
The loop is deliberately slow --- it adapts over many coordination
events, not within one --- which keeps it cleanly separated from the
per-event coordination loop of the Marketplace.

\paragraph{Candidate mechanisms.}
Two mechanism families are natural candidates and are recorded here as
the substance of the future-work proposal. The first is \emph{Bayesian
weight adaptation}: treat each agent's desire weights as a posterior
distribution updated from observed outcomes, so that the weights track
the plant's operating regime with a principled account of uncertainty.
The second addresses a hazard the first does not. If the weights are
simply re-fitted to the current regime, the agents may \emph{forget}
how to coordinate well in a regime that recurs --- catastrophic
forgetting. Elastic Weight Consolidation \cite{kirkpatrick2017} is the
reference mechanism for exactly this problem: it slows adaptation of the
weights that were important for previously seen regimes, so that the
agents accumulate competence across regimes rather than overwriting it.
A School layer built on Elastic Weight Consolidation would let \NEXUS{}
adapt to a changing plant without losing coordination skill it had
already acquired.

\paragraph{Why it is deferred.}
The School layer is deferred for scope reasons, not because it is
ill-defined. The three implemented layers already constitute a complete,
testable reference architecture for the cross-functional coordination
problem; continual adaptation is a second-order improvement on top of a
working coordinator, and folding it into the present scope would broaden
the empirical work beyond what the thesis timeline allows. Deferring it
is consistent with the treatment of the School-like adaptation slot in
the holonic-architecture lineage, where the adaptive element is
typically named in the architecture and elaborated separately. The
detailed future-work treatment is given in
Chapter~\ref{ch:conclusions}.

% ------------------------------------------------------------
\section{ISA-95 Mapping and Privacy Preservation}
\label{sec:arch-isa95}

This section consolidates two cross-cutting properties of the
architecture that the preceding sections established piecewise: the
mapping of \NEXUS{} onto the ISA-95 Level~3 functional model, and the
privacy-preservation property that the layered, bid-mediated design
delivers.

\subsection{Mapping onto ISA-95 Level~3}
\label{subsec:isa95-map}

ISA-95 Part~1 places Manufacturing Operations Management at Level~3 of
the functional hierarchy and decomposes it into operations-management
activities for production, maintenance, quality and inventory
\cite{isa95part1,sprock2019isa95}. \NEXUS{} maps onto this model
directly, and the directness of the mapping is itself a design
argument: an architecture whose components are in one-to-one
correspondence with the standard's functions can be explained to, and
owned by, the people who already operate those functions.

\begin{table}[ht]
\centering
\caption{Mapping of ISA-95 Level~3 operations-management functions onto
\NEXUS{} components. Each function becomes one Worker agent owning one
contract issue; the cross-function coordination the standard leaves
informal becomes the explicit task of the Marketplace layer.}
\label{tab:isa95}
\small
\begin{tabularx}{\textwidth}{@{}l l l X@{}}
\toprule
\textbf{ISA-95 L3 function} & \textbf{\NEXUS{} agent} & \textbf{Owned issue}
  & \textbf{Decision authority} \\
\midrule
Production operations management & \PSA{} & \texttt{production\_rate}
  & Sets the production rate; subordinate voice on the other three issues. \\
\addlinespace
Maintenance operations management & \PMA{} & \texttt{maint\_window}
  & Sets the maintenance window; subordinate voice elsewhere. \\
\addlinespace
Quality operations management & \QCA{} & \texttt{quality\_action}
  & Sets the quality hold; subordinate voice elsewhere. \\
\addlinespace
Inventory / supply operations management & \SCA{} & \texttt{supply\_horizon}
  & Sets the supply horizon; subordinate voice elsewhere. \\
\addlinespace
Cross-function coordination (informal in ISA-95) & Marketplace layer
  & --- & Resolves the joint decision as one composite contract. \\
\addlinespace
Operations safety / governance & Fence layer & ---
  & Audits each contract against pre-declared safety rules. \\
\bottomrule
\end{tabularx}
\end{table}

Table~\ref{tab:isa95} records the mapping. Each of the four Level~3
functions becomes one Worker agent that owns one issue of the composite
contract --- the operational realisation of the ``cross-functional
decision authority'' notion: an agent has decision authority over its
home issue and a subordinate, damped voice over the others
(Section~\ref{subsec:mono-dim}). The element ISA-95 leaves implicit ---
the coordination \emph{between} the four functions --- is exactly what
\NEXUS{} makes explicit and assigns to the Marketplace layer. The
standard acknowledges that the functions must be reconciled but does
not prescribe a mechanism; \NEXUS{} contributes one. The Fence layer,
finally, occupies the operations-safety slot, auditing the coordinated
decision against rules that sit outside any single function's remit.

\subsection{Privacy preservation}
\label{sec:arch-privacy}
\label{subsec:privacy}

The privacy property of \NEXUS{} can now be stated precisely, because
all the pieces have been introduced. By construction of the Worker and
Marketplace layers:

\begin{quote}
\itshape
For each coordination event, an agent discloses only its bid set --- a
bounded number of (utility, region) pairs. It never discloses its
beliefs, its constraint set, its constraint weights, or its desire
weights. The Mediator's entire information set is the union of the bid
sets disclosed for that event.
\end{quote}

The disclosure is bounded in a strong sense: the number of regions an
agent publishes per event is capped by a protocol parameter and is
independent of the size of the agent's underlying belief and constraint
model. An agent with a rich internal model and an agent with a sparse
one publish bid sets of the same bounded shape. The Mediator therefore
never learns an agent's utility \emph{function}; it learns only a
bounded sample of \emph{values} of that function, at regions the agent
itself chose to nominate.

Table~\ref{tab:disclosure} places this property against the two
reference points. A centralised optimiser such as NSGA-III
\cite{deb2014nsga3} --- used as the centralised comparator in
Chapter~\ref{ch:results} --- requires the entire joint problem: every
function's cost structure, in full. Distributed constraint optimisation
\cite{modi2005adopt,fioretto2018dcop} distributes the computation but
still requires each agent to publish complete cost tables over the
shared variables. \NEXUS{} requires neither: it discloses bid summaries
only. It occupies a distinct point in the disclosure--coordination
trade-off space --- coordinated joint decisions under bounded
disclosure --- and that point is the architectural contribution the
privacy property names.

\begin{table}[ht]
\centering
\caption{Information disclosure under three coordination regimes.
\NEXUS{} produces a coordinated joint decision while disclosing only
bounded bid summaries.}
\label{tab:disclosure}
\small
\begin{tabularx}{\textwidth}{@{}l X c@{}}
\toprule
\textbf{Regime} & \textbf{What each function must disclose} & \textbf{Joint decision} \\
\midrule
Centralised optimiser (NSGA-III) &
  Full cost structure, state, and models --- the entire joint problem. & yes \\
\addlinespace
DCOP &
  Complete cost tables over all shared variables. & yes \\
\addlinespace
\textbf{\NEXUS{}} &
  Bid summaries only: a bounded set of (utility, region) pairs per
  event. Beliefs, constraints, weights stay private. & yes \\
\addlinespace
Siloed operation &
  Nothing --- but no joint decision is produced. & no \\
\bottomrule
\end{tabularx}
\end{table}

Two qualifications keep the claim honest. First, the privacy of
\NEXUS{} is \emph{organisational}, not \emph{cryptographic}. The claim
is that the architecture never \emph{requires} an agent to publish its
model --- not that a determined adversary observing many events could
learn nothing. An observer who collected an agent's bid sets across a
long run could in principle reconstruct a coarse approximation of the
agent's utility function from the nominated regions and their annotated
values. Quantitative, differential-privacy-style guarantees over bid
disclosure are out of scope and are noted as future work in
Chapter~\ref{ch:conclusions}. Second, the bound is on disclosure
\emph{per event}: it is the shape of a single event's disclosure that is
independent of model size, and the cumulative disclosure over a run
grows with the number of events, as the first qualification reflects.
With those qualifications stated, the property stands: \NEXUS{}
coordinates four functions into a joint decision while requiring each
to disclose only bounded bid summaries, and that is what no centralised
or DCOP-style alternative achieves.

% ------------------------------------------------------------
\section{Implementation Architecture}
\label{sec:arch-impl}

The architecture described in this chapter is realised as a Python
codebase; this section gives an overview of its structure and forward
references to Chapter~\ref{ch:implementation}, where the implementation
is described in full.

The codebase mirrors the four-layer decomposition. The Worker layer
lives in an \texttt{agents/} package: a base BDI agent encapsulates the
shared belief-update and bid-generation pipeline, and four subclasses
--- one per agent --- contribute the specialised belief update and
constraint derivation. The Marketplace layer lives in a
\texttt{marketplace/} package: the Bid Bus, the Mediator with the
branch-cut A* search of Section~\ref{subsec:mediator}, and the Contract
Record. The Fence layer lives in a \texttt{fence/} package: the rule
suite and the three-tier classifier. The composite contract, the bid,
and the constraint --- the data structures shared across layers --- live
in a \texttt{core/} package, and the discrete-event plant model lives in
an \texttt{environment/} package. A \texttt{baselines/} package holds
the comparison mechanisms used in Chapter~\ref{ch:results}: the siloed
(uncoordinated) baseline, the rule-based coordination baseline, and a
single-winner contract-net mechanism. The one-to-one correspondence
between layers and packages is deliberate: the architectural boundaries
of this chapter are the module boundaries of the code, so a reader of
the implementation can navigate it with the four-layer diagram of
Figure~\ref{fig:arch-layers} in hand.

Two properties of the implementation are worth previewing here because
they bear on the credibility of the architecture as a whole.

The first is \emph{test coverage}. The implementation carries an
automated regression suite of more than $240$ tests spanning every
layer: the core data structures, the mediator search, the environment,
the Fence classifier, the four agents, the baseline mechanisms, and
end-to-end integration runs. The suite is what lets the architecture be
modified with confidence --- the admissible-heuristic correction of
Section~\ref{subsec:mediator}, for instance, was made and verified
against it --- and it is described in detail in
Chapter~\ref{ch:implementation}.

The second is \emph{reproducibility}. A coordination event in \NEXUS{}
is fully deterministic: given the same conflict announcement and the
same agent state, the agents generate the same bids --- the bid
generation defaults to deterministic greedy search
(Section~\ref{subsec:bid-gen}) --- the Mediator runs the same search,
and the deterministic tie-break selects the same contract. Determinism
holds across processes, because the agent-side random sampling is keyed
by a stable hash rather than by a process-dependent seed. The practical
consequence, quantified in Chapter~\ref{ch:implementation}, is that a
whole experimental batch reproduces byte-for-byte on re-execution. For a
reference architecture this is not a mere convenience: it means every
empirical claim in Chapter~\ref{ch:results} rests on a computation that
any reader can re-run and obtain identically.

\subsection*{Summary}

This chapter has presented \NEXUS{} as a four-layer reference
architecture for cross-functional coordination at ISA-95 Level~3. The
Worker layer represents each of the four Level~3 functions as a BDI
agent with mono-dimensional authority over one issue of a composite
contract. The Marketplace layer resolves each cross-functional conflict
into a single composite contract by an admissible-heuristic search that
maximises joint welfare over the agents' published bids, with a
global-optimality guarantee and a consequent Pareto-optimality
property. The Fence layer audits every contract against a suite of
pre-declared safety rules and classifies it into a three-tier severity
scheme, integrated structurally as a log-only audit layer with no
coupling into the negotiation protocol. The School layer, deferred as
future work, provides the architectural slot for continual adaptation
of the agents' preferences. Across all four layers the architecture
holds to bounded information disclosure: each function reveals only
bid summaries, never its internal model, so \NEXUS{} produces a
coordinated joint decision without the full-disclosure requirement that
makes centralised optimisation infeasible in practice. The next chapter
turns from the architecture to its implementation.
